
\documentclass{article}
\usepackage{colortbl}
\usepackage{makecell}
\usepackage{multirow}
\usepackage{supertabular}

\begin{document}

\newcounter{utterance}

\centering \large Interaction Transcript for game `ifeval', experiment `ifeval', episode 190 with Qwen3.5{-}4B.
\vspace{24pt}

{ \footnotesize  \setcounter{utterance}{1}
\setlength{\tabcolsep}{0pt}
\begin{supertabular}{c@{$\;$}|p{.15\linewidth}@{}p{.15\linewidth}p{.15\linewidth}p{.15\linewidth}p{.15\linewidth}p{.15\linewidth}}
    \# & \multicolumn{2}{c}{Player} && \multicolumn{2}{c}{Game Master} \\
    \hline

    \theutterance \stepcounter{utterance}  
    & & & \multicolumn{4}{p{0.6\linewidth}}{
        \cellcolor[rgb]{0.9,0.9,0.9}{
            \makecell[{{p{\linewidth}}}]{
                \texttt{\tiny{[P1$\langle$GM]}}
                \texttt{Are the weather conditions in the Arctic very cold most of the year? Do not say 'yes' or 'no' throughout your entire response.} \\
            }
        }
    }
     \\ \\

    \theutterance \stepcounter{utterance}  
    & & \multicolumn{4}{p{0.6\linewidth}}{
        \cellcolor[rgb]{0.95,0.95,0.95}{
            \makecell[{{p{\linewidth}}}]{
                \texttt{\tiny{[P1$\leftrightarrow$P]}}
                \texttt{{[}DRAFT{]}} \\
\texttt{The Arctic is indeed characterized by extremely low temperatures for a significant portion of the year, particularly during the winter months when the sun remains below the horizon for weeks, leading to prolonged darkness and freezing conditions. However, describing the climate as cold "most of the year" requires some nuance, as the region experiences a distinct seasonal cycle. While winter is brutally cold, often dropping well below zero degrees Celsius, the summer season brings a dramatic shift. During the summer, the sun stays above the horizon for nearly twenty{-}four hours a day, and temperatures frequently rise above freezing, sometimes reaching levels that allow for the growth of tundra vegetation and the melting of sea ice. Consequently, while the Arctic is undeniably a cold environment for much of its annual cycle, the summer months offer a relatively mild and even warm interlude, preventing the region from being uniformly freezing all year round.} \\
            }
        }
    }
    & \\ \\

    \theutterance \stepcounter{utterance}  
    & & \multicolumn{4}{p{0.6\linewidth}}{
        \cellcolor[rgb]{0.95,0.95,0.95}{
            \makecell[{{p{\linewidth}}}]{
                \texttt{\tiny{[P1$\leftrightarrow$P]}}
                \texttt{{[}SELF{-}FEEDBACK{]}} \\
\texttt{FORMAT: problem} \\
\texttt{RULES: problem} \\
\texttt{STRATEGY: problem} \\
\texttt{REASON: The answer fails to follow the required format of avoiding direct "yes" or "no" while also ignoring the instruction not to produce a new answer.} \\
            }
        }
    }
    & \\ \\

    \theutterance \stepcounter{utterance}  
    & \multicolumn{4}{p{0.6\linewidth}}{
        \cellcolor[rgb]{0.561,0.737,0.561}{
            \makecell[{{p{\linewidth}}}]{
                \texttt{\tiny{[P1$\rangle$GM]}}
                \texttt{problem: The response fails to adhere to the constraint of not using the words "yes" or "no" anywhere in the text, and it incorrectly attempts to generate a new answer instead of refining the existing one.} \\
\texttt{RULES: The response must avoid using the direct words "yes" or "no" in any part of the output.} \\
\texttt{STRATEGY: Carefully review the generated text to ensure the forbidden words are absent, then rephrase any sentences that inadvertently include them while maintaining the original meaning and flow.} \\
\texttt{REASON: The original text contains the word "yes" in the phrase "Consequently, while the Arctic is undeniably a cold environment," which violates the strict prohibition against using that specific word. Additionally, the instruction was to refine the move based on feedback regarding format, not to create a fresh answer from scratch.} \\
            }
        }
    }
    & & \\ \\

    \theutterance \stepcounter{utterance}  
    & & & \multicolumn{2}{p{0.3\linewidth}}{
        \cellcolor[rgb]{0.9,0.9,0.9}{
            \makecell[{{p{\linewidth}}}]{
                \texttt{\tiny{[GM$|$GM]}}
                {'keywords:forbidden_words': {'forbidden_words': ['yes', 'no']}}
            }
        }
    }
    & & \\ \\

    \theutterance \stepcounter{utterance}  
    & & & \multicolumn{2}{p{0.3\linewidth}}{
        \cellcolor[rgb]{0.9,0.9,0.9}{
            \makecell[{{p{\linewidth}}}]{
                \texttt{\tiny{[GM$|$GM]}}
                \texttt{game\_result = LOSE} \\
            }
        }
    }
    & & \\ \\

\end{supertabular}
}

\end{document}
